% 두 뱅크가 STROBE·JUDGE·CLEAR를 만드는 시간표 (그림 fig:hw-timing-banks). 가로 = ms, 두 그림이 같은 눈금.
% (가) 릴레이식 순서기: 링 계수기 Q0--Q9, 한 칸 6.3 ms, 주기 63 ms (hardware/relay/sequencer/README.md).
% (나) 기어식: 60 rpm, 한 위치 38.46 ms, STROBE 접점 시작 = 0. 슬롯 = 커넥터 면
%      (docs/decisions/0003-interfaces.md 3.4절, hardware/gear/layout.md 3.8절).
\begin{tikzpicture}[x=1.95mm,y=1cm,
    lab/.style={font=\footnotesize,anchor=east},
    nt/.style={font=\scriptsize,text=cgray},
    bar/.style={draw=#1,fill=#1!25,line width=0.5pt},
    face/.style={draw=#1,dashed,line width=0.5pt}]
  \def\tbar#1#2#3#4{\draw[bar=#4] (#2,{#1-0.17}) rectangle (#3,{#1+0.17});}
  % ================= (가) 릴레이식 순서기 =================
  \node[font=\footnotesize\bfseries,anchor=west] at (-10,5.2) {(가) 릴레이식: 순서기의 링 계수기 (한 칸 6.3 ms, 한 위치 63 ms)};
  \foreach \q in {0,...,9} {
    \draw[draw=cgray,fill=white] ({\q*6.3},4.45) rectangle ({\q*6.3+6.3},4.85);
    \node[font=\scriptsize] at ({\q*6.3+3.15},4.65) {Q\q};
  }
  \node[lab] at (-0.6,4.65) {링의 칸};
  \tbar{3.95}{0}{12.6}{cblue}       \node[lab] at (-0.6,3.95) {\code{STROBE}};
  \tbar{3.45}{18.9}{31.5}{cgreen}   \node[lab] at (-0.6,3.45) {\code{JUDGE}};
  \tbar{2.95}{37.8}{56.7}{corange}  \node[lab] at (-0.6,2.95) {\code{CLEAR}};
  \tbar{2.45}{44.1}{56.7}{cpurple}  \node[lab] at (-0.6,2.45) {\code{CK}};
  \draw[cpurple,->] (37.8,2.15) -- (37.8,2.3);
  \node[nt,anchor=north] at (37.8,2.15) {Q6에서 \code{CK}를 낼지 정함};
  \node[nt,anchor=west] at (57.2,2.45) {전진};
  \draw[cgray,dashed] (63,2.2) -- (63,4.9);
  \node[nt,anchor=south] at (63,4.9) {다음 위치};
  % ================= (나) 기어식 슬롯 =================
  \begin{scope}[yshift=-3.6cm,xshift=5.85mm]
  \node[font=\footnotesize\bfseries,anchor=west] at (-13,4.75) {(나) 기어식: 슬롯 원판 (빠른 드럼 60 rpm, 한 위치 38.46 ms)};
  \draw[draw=cgray,fill=clight] (-3,4.0) rectangle (9,4.3);
  \node[font=\scriptsize] at (3,4.15) {패드 위};
  \node[lab] at (-3.6,4.15) {드럼 접점};
  \draw[cred,line width=0.7pt] (16.2,3.85) -- (16.2,4.45);
  \node[font=\scriptsize,text=cred,anchor=west] at (16.4,4.15) {접점이 패드를 떠남 (약 +16.2)};
  \tbar{3.55}{-2.50}{4.52}{cblue}     \draw[face=cblue] (0,3.38) rectangle (6,3.72);
  \tbar{3.55}{35.96}{38.46}{cblue}
  \node[lab] at (-3.6,3.55) {\code{STROBE}};
  \tbar{3.05}{8.56}{16.63}{cgreen}    \draw[face=cgreen] (11.06,2.88) rectangle (18.13,3.22);
  \node[lab] at (-3.6,3.05) {\code{JUDGE}};
  \tbar{2.55}{18.65}{31.73}{corange}  \draw[face=corange] (21.15,2.38) rectangle (33.23,2.72);
  \node[lab] at (-3.6,2.55) {\code{CLEAR}};
  \draw[->,cgray] (-3,1.95) -- (40,1.95) node[right,font=\scriptsize] {ms};
  \foreach \t in {0,10,20,30} { \draw[cgray] (\t,1.9) -- (\t,2.0); \node[nt,anchor=north] at (\t,1.9) {\t}; }
  \draw[cgray,dashed] (38.46,1.95) -- (38.46,4.4);
  \end{scope}
  % (가)의 시간축
  \draw[->,cgray] (0,1.75) -- (66,1.75) node[right,font=\scriptsize] {ms};
  \foreach \t in {0,10,...,60} { \draw[cgray] (\t,1.7) -- (\t,1.8); \node[nt,anchor=north] at (\t,1.7) {\t}; }
\end{tikzpicture}
